\section{Training strategies}\label{sec:training}

\subsection{Behaviour cloning (the teammates' baselines)}
Procedure (\pth{bev_policy/collect_demos.py}, \pth{train_bc.py}): the rule-based driver (stock intelligent driver model) drives a small number of episodes; at every step the observation (grid and six scalars from the teammates' perception chain) and the driver's action are stored.
Only episodes that reach the destination are kept (arrival filter). The network is trained by regression on the action (mean squared error), 60 epochs maximum, batch size 64, learning rate $3\cdot10^{-4}$, early stopping with patience 10.
The two repository baselines use \textbf{five demonstrations each}: one set without other vehicles (\code{repo_bc_clear}), one with sparse traffic (\code{repo_bc_sparse}).

\textbf{Why this is weak.} The network is trained only on states that the expert visits. Once the learner makes a small error it reaches states the expert never visited and has no labelled example of how to recover (covariate shift; Ross et al., 2011).
A concrete symptom was documented early: the cloned policy creeps to 0.8--0.96\,m from the lane centre by step 80--100 (1.58\,m for seed 1001), a systematic lateral bias. With only five demonstrations the policy also never sees most traffic situations.

\subsection{Dataset aggregation (the main training strategy)}
Dataset aggregation fixes the covariate shift by letting the \emph{learner} choose which states are visited and asking the \emph{teacher} what it would have done there. Implementation: \pth{bev_rl/adaptive/dagger_base.py} (and the variants listed below).

\begin{lstlisting}[caption={Dataset aggregation as implemented (training seeds 2000--2999, validation seeds 3100--3109).},label={lst:dagger}]
data = empty;  student = fresh network (random initialisation)
for round t = 0 .. 5:
    beta = 1.0 if t == 0 else 0.5 ** t          # probability that the TEACHER's action is executed
    for each traffic density d in {0, 0.05, 0.1, 0.2, 0.3}:
        run 8 episodes on randomly drawn training seeds:
            at every step:  record (grid, scalars, teacher_action)      # teacher sees the true state
                            execute teacher_action with probability beta, else student_action
                            add Gaussian noise (std 0.1) to the executed action   # widen the visited states
    data = data + new records
    student = train(student, data):  Adam, learning rate 1e-3, batch 128, mean squared error,
                                     8 epochs (24 epochs in round 0), on ALL data so far
    validate: 10 episodes (seeds 3100-3109) at densities 0.05 and 0.2, no teacher
    keep as 'best' the round with the highest mean success (ties: fewer collisions + off-road)
\end{lstlisting}
The aggregated data set holds about 67\,000 labelled steps after round 5 for the unmodified-grid student and 73\,032 for the headline model. One training seed per model was run (seed 0); differences between rounds and between models of 0.1 in validation success correspond to a single episode of ten and must not be over-read (see Figure~\ref{fig:dagger}).

Three properties matter for interpreting the results:
\begin{itemize}
\item The teacher is a \emph{rule-based} controller that sees the exact simulator state; the student only sees the grid and a few scalars. The student can therefore at best match the teacher, and wherever the information is not in the student's input (for example a vehicle 30\,m away) it cannot imitate the teacher's reaction (Section~\ref{sec:failure}).
\item Model selection uses only validation seeds; the test seeds (1000--1049, 5000--5199) are never touched during training or selection.
\item In round 0 the teacher drives ($\beta=1$) and the student learns from the teacher's own distribution; the later rounds are what supply recovery behaviour.
\end{itemize}

\begin{figure}[H]\centering
\includegraphics[width=\textwidth]{figures/fig_dataset_aggregation_validation_curves.pdf}
\caption{Validation success per dataset-aggregation round for five input designs (10 validation episodes each; one training seed). In sparse traffic (left) all designs except ``map overlay, no scalars'' are at 1.0 in almost every round, so validation cannot discriminate between them; in dense traffic (right) values of 0.0--0.5 differ by at most five episodes of ten.
The ``map overlay, no scalars'' design is clearly worse in both panels. Lines that coincide are drawn on top of each other.}
\label{fig:dagger}
\end{figure}

\subsection{Variants of the dataset-aggregation recipe}
All variants keep the optimiser, loss, action noise and network; they change one thing at a time.
\begin{table}[H]\centering\small
\caption{Variants of the dataset-aggregation recipe. ``Teacher'' is the source of labels; ``perception'' is what the student sees.}\label{tab:variants}
\begin{tabularx}{\textwidth}{@{}>{\raggedright\arraybackslash}p{3.7cm}>{\raggedright\arraybackslash}p{3.6cm}>{\raggedright\arraybackslash}p{2.8cm}X@{}}\toprule
Variant & Perception & Teacher & Specific settings\\\midrule
Input ablations (\code{base_a}, \code{ov_both}, \code{ov_bev_only}, \code{arch_*}) & Ideal detector; unmodified or overlay grid; all, no or a subset of scalars (inputs hidden by constants) & Stock model & Densities 0--0.3; validation densities 0.05 and 0.2.\\
Headline model \code{ov_split} & Ideal detector; overlay grid; scalars speed, progress, goal & Stock model & As above.\\
Safe-teacher student & As headline & Collision-aware teacher & Densities 0, 0.05, 0.1; validation 0.05; 7 rounds (0--6).\\
Caution study \code{caut_a}, \code{ctrl_nocaut} & Detector dropout at measured retention per weather; overlay grid & Caution-aware stock model (caution 0: 30\,km/h, 10\,m gap; caution 1: 20\,km/h, 25\,m gap) & Caution value per episode: half uniform in $[0,1]$, half derived from weather as $1-\text{retention}$; control network has the input fixed to 0.\\
Degraded-perception student \code{deg_base} & Calibrated degradation (missed detections and range errors, Section~\ref{sec:perception}); overlay grid & Collision-aware teacher on ground truth & Training weather mix: clear (twice), fog 0.8 and 1.0, rain 0.3, 0.6, 1.0 (twice); densities 0.05, 0.05, 0.1, 0; 5 rounds; 6 epochs.\\
Pose-error students (\code{pose_*}) & Ideal detector; overlay drawn from a \emph{wrong} car pose & Stock model on the true state & Per episode an error level $L$ ($30\%$ exactly 0, 70\% uniform): lateral error $\pm1.5L$\,m, heading error $\pm0.15L$\,rad; conditioned variants receive $L$; two-cue variants additionally receive noisy lateral and heading scalars.\\
Pixel policy \code{vcam} & Front-camera image & Collision-aware teacher & Behaviour cloning on clean images with action noise (0, 0.1, 0.25), then three dataset-aggregation rounds.\\\bottomrule
\end{tabularx}
\end{table}

\subsection{Reinforcement learning}
Four reinforcement-learning attempts are described in Section~\ref{sec:rl}. Common to all of them is the algorithm (proximal policy optimisation from Stable-Baselines3: 256 steps per update per environment, batch 128, 5 epochs, discount 0.99, GAE parameter 0.95, clip range 0.1, no entropy bonus, target KL 0.02; the first attempt used a plain multilayer-perceptron policy with 512 steps per update, batch 512, 10 epochs, clip range 0.2 and target KL 0.03),
and the idea of starting from an imitation-trained network instead of from random weights, because a random start settles in a stand-still optimum (zero speed is the safest behaviour).
